\subsection*{1. Envoi}
The opening locates the program within Grothendieck's application to the
CNRS and explains the work he intends to carry forward. He refers to the
1972 account of his ``Titres et Travaux'', with its list of publications,
and to two boxes of subsequent manuscript notes that have passed through
successive stages of ``décantation''. His account of ``rédaction'' already
qualifies the word: on p.~1 he describes the projected writing as the
further development, clarification, and deepening of ideas, with the
unexpected turns that accompany discovery. The ``Réflexions Mathématiques'',
begun in February 1983 according to this passage, are to make that movement
of thought visible. The sentence crossing pp.~1--2 gives the complementary
operation: an inadequate image is exposed, and the objects previously
encountered askew acquire their mutual coherence. The final paragraph
assigns the ensuing sketch both a retrospective role, covering ten or
twelve years of reflection, and a prospective role as a program for the
coming years. Its stated readers include the Comité National and other
colleagues, former students, and friends. The 1972 report and the manuscript
``Réflexions Mathématiques'' are named here as the author's own points of
reference; the checked passage supplies their role in this application.
\S7 returns to the projected \emph{Réflexions Mathématiques}, while
\S\S9--10 return to the teaching appraisal and the CNRS application.
(\href{https://grothendieck.umontpellier.fr/web/viewer.html?file=../119.pdf\#page=4}{Copy 1, p.~1, site 3};
\href{https://grothendieck.umontpellier.fr/web/viewer.html?file=../119.pdf\#page=5}{p.~2, site 4}.)

